\chapter{Hardware, alimentación y montaje}
\label{chap:hardware}

\section{Inventario del montaje}

Antes de empezar, anota los datos que puedas comprobar de cada componente.
No completes una referencia a partir de una fotografía o por parecido con otro modelo.
Si algún dato, como la revisión de una placa o su tensión de alimentación, no puede
confirmarse con la etiqueta, la factura o la documentación del fabricante, indícalo como
\emph{no verificado}.

\begin{table}[H]
\centering
\caption{Datos básicos del montaje.}
\label{tab:inventario-lector}
\begin{tabularx}{\textwidth}{@{}L{0.28\textwidth}L{0.31\textwidth}Y@{}}
\toprule
\textbf{Elemento} & \textbf{Datos} & \textbf{Fuente} \\
\midrule
Host & Fabricante, modelo, RAM y revisión & Equipo, firmware o factura \\
Sistema & Distribución, versión y arquitectura & \code{/etc/os-release}, \code{uname -m} \\
Kernel & Versión completa & \code{uname -r} \\
Tarjeta Akida & Modelo y revisión & Etiqueta y documentación de BrainChip \\
Adaptador & Fabricante, modelo y tipo de enlace & Serigrafía y documentación del fabricante \\
Cable & Tipo y orientación & Marcado y documentación del adaptador \\
Fuente del host & Tensión y potencia & Etiqueta de la fuente \\
Fuente de la tarjeta o carrier & Tensión, polaridad y conector & Etiqueta y documentación del fabricante \\
\bottomrule
\end{tabularx}
\end{table}

Los conectores M.2, Mini PCIe, PCIe convencional y el conector FFC de Raspberry Pi 5
no son equivalentes. Comprueba siempre el tipo de conexión y la orientación del cable
en la documentación del equipo utilizado.

\section{Topologías cubiertas}

\subsection{Host Linux con ranura PCIe}

En un PC Linux, la tarjeta se instala en una ranura PCIe compatible y se alimenta según
las indicaciones de su documentación. El equipo debe estar apagado y desconectado antes
de insertar o retirar la tarjeta.

Después del montaje, continúa con el \cref{chap:host-driver}. Si \code{lspci}
no detecta un nuevo dispositivo PCIe, el problema debe resolverse antes de continuar
con la configuración de Python.

\subsection{Raspberry Pi 5 con adaptador FFC a PCIe}

Raspberry Pi 5 dispone de una interfaz PCIe Gen~2~$\times$1 accesible mediante un
conector FFC \cite{rpi_pcie}. Para utilizar una tarjeta PCIe convencional es necesario
un adaptador que lleve este enlace hasta una ranura compatible.

Aunque el adaptador disponga de una ranura físicamente más larga, el enlace proporcionado
por Raspberry Pi 5 sigue siendo $\times$1.

El montaje del TFG se muestra en la \cref{fig:montaje-aislado}.

\begin{figure}[H]
  \centering
  \includegraphics[width=0.92\textwidth]{hardware-carrier-detail-isolated.png}
  \caption[Montaje Raspberry Pi 5, carrier y AKD1000]{Montaje formado por una Raspberry Pi 5,
  un adaptador PCIe y la tarjeta BrainChip AKD1000 utilizado durante el TFG.}
  \label{fig:montaje-aislado}
\end{figure}

\section{Montaje}

Si la documentación del fabricante indica una secuencia de conexión diferente,
sigue ese procedimiento.

\begin{enumerate}
  \item Apaga el host y desconecta las fuentes de alimentación.

  \item Comprueba la orientación del cable FFC y la tensión de alimentación del
        adaptador o de la tarjeta.

  \item Conecta el cable FFC al adaptador respetando la orientación indicada.

  \item Inserta la AKD1000 completamente en la ranura PCIe del carrier.

  \item Revisa las conexiones antes de alimentar de nuevo el sistema.
\end{enumerate}

\begin{stopbox}{Si aparece algún problema eléctrico}
Desconecta la alimentación si aparece humo, olor anormal, una temperatura excesiva
o cualquier otro comportamiento que indique un posible problema de alimentación
o montaje.
\end{stopbox}

\section{Carrier Waveshare utilizado en el TFG}

En el montaje del TFG se utilizó un adaptador Waveshare conectado a la Raspberry Pi 5
mediante un cable FFC de 50~mm. En la placa puede leerse \code{DC 12V} y
\emph{PCIe TO PCIe x1 Board (C)}.

La documentación actual de Waveshare describe este modelo con una ranura compatible
mecánicamente con tarjetas PCIe $\times$1, $\times$4, $\times$8 y $\times$16
\cite{waveshare_board_c}. El enlace con Raspberry Pi 5 continúa siendo, en cualquier caso,
PCIe $\times$1.

No se conserva el número de artículo exacto del adaptador utilizado. La identificación
disponible del montaje se recoge en el \cref{tab:componentes-tfg}.

\begin{figure}[H]
  \centering
  \includegraphics[width=0.78\textwidth]{hardware-powered-isolated.png}
  \caption[Montaje alimentado]{Montaje utilizado durante una de las sesiones del TFG.}
  \label{fig:montaje-alimentado}
\end{figure}

\section{Comprobación antes del encendido}

Antes de alimentar el montaje, comprueba que el cable FFC está correctamente orientado,
que la tarjeta está bien insertada y que las tensiones de alimentación coinciden con las
indicadas para los componentes utilizados.

\begin{stopbox}{Condición de parada H0}
Si no puedes confirmar la tensión, la polaridad o la orientación de alguna conexión,
no alimentes el montaje hasta comprobarlo en la documentación correspondiente.
\end{stopbox}